% ECONOMICS_PROBLEM: {"id": "C78-464", "title": "Rank-maximality probability under uncertain preferences", "jel_code": "C78", "source_id": "OP-0548"}
% FIRST_POSTED: 2026-10-10
% ATTEMPT_STATUS: solved
% ENTRY_KIND: research_attempt
\documentclass[10pt]{article}
\usepackage[a4paper,margin=24mm]{geometry}
\usepackage{amsmath,amssymb,amsthm}
\usepackage[T1]{fontenc}
\usepackage{lmodern}
\usepackage{microtype}
\usepackage[colorlinks=true,linkcolor=blue,citecolor=blue,urlcolor=blue]{hyperref}
\newtheorem{theorem}{Theorem}
\newtheorem{lemma}{Lemma}
\newcommand{\IS}{\mathcal I}
\newcommand{\FP}{\mathsf{FP}}
\newcommand{\PP}{\mathsf{PP}}
\newcommand{\sharpP}{\mathsf{\#P}}
\title{Rank-maximality probabilities under uncertain preferences}
\author{Ji Ho Bae}
\date{10 October 2026}
\begin{document}
\maketitle

\paragraph{Model, authorship and original output.}
Author: Ji Ho Bae. Prepared with GPT Codex assistance; the precise serving
revision was not exposed. The final generated mathematical argument is
retained at the following immutable source link. This catalogue submission
preserves that argument, adding catalogue metadata and provenance; the
rank-maximality title layout is simplified for browser rendering.
\href{https://github.com/jbaelaw/rank-maximality-hand-proof-20261010/blob/e0d34b34f26033d6cef32493a65abca1480a37b2/paper/Proof.tex}{Complete original generated TeX}.
\href{https://github.com/jbaelaw/rank-maximality-hand-proof-20261010/blob/e0d34b34f26033d6cef32493a65abca1480a37b2/paper/Proof.pdf}{Hand-proof PDF}.
The preparation included exact target and primary-source checks, adversarial
mathematical reviews, and finite semantic checks recorded in the public
package. This is a complete hand-proof claim under polynomial-time Turing
reductions, submitted for independent expert review. It does not claim
community acceptance or an end-to-end Lean formalization.

\begin{abstract}
For a specified house allocation, we classify the probability of rank
maximality under uncertain strict preferences. Both independent rational
lotteries and independent uniform tie refinements give an
$\FP^{\sharpP}$-complete exact evaluation problem and a
$\PP$-complete threshold problem under polynomial-time Turing reductions.
The lower bound uses only two-house lists: one exact probability reveals
the independent-set count as a sum of integer digits. Explicit joint
lotteries admit polynomial-time evaluation and comparison.
\end{abstract}

Applicants rank their acceptable houses; houses have no preferences.
Acceptability is fixed. A feasible matching may leave applicants or houses
unmatched. Its rank vector counts the assigned rank-one edges, then the
rank-two edges, and so on, padded to the largest acceptable-list length.
It is \emph{rank-maximal} when this vector is lexicographically maximal
over all feasible matchings. Cardinality is not maximized beforehand.
For a fixed feasible matching $M$, write $\eta(M)$ for the probability
of this event. The survey \cite[\S3.1.2]{survey} proposes studying
rank-maximal matchings under uncertainty; catalogue entry C78-464
specifies the probability-classification problem considered here.

\begin{theorem}\label{thm:main}
For each of the following independent uncertainty models, exact rational
evaluation of $\eta(M)$ is $\FP^{\sharpP}$-complete and deciding
$\eta(M)\geq q$ for a rational input $q$ is $\PP$-complete, under
polynomial-time Turing reductions:
\begin{enumerate}
\item each applicant has an explicitly listed rational lottery over strict
orders;
\item each applicant has a weak order, independently and uniformly refined
within every tie.
\end{enumerate}
Hardness holds with at most two acceptable houses per applicant and,
respectively, fair two-outcome nontrivial lotteries or nontrivial ties
of size two. For an
explicitly listed lottery over whole profiles, both tasks are polynomial-time.
\end{theorem}

\section*{The two-choice construction}
Let $G=(V,E)$ be a finite simple graph, with $n$ vertices and $m$ edges,
and let $t\geq1$. For each $v\in V$, create houses $h_v,k_v$, a variable
applicant $a_v$ accepting $h_v,k_v$, and a guard $b_v$ accepting only $k_v$.
For each edge $\{u,v\}$, create $t$ additional applicants accepting
$h_u,h_v$. Every two-house list is a single tie of size two, refined
independently and fairly. The same distribution is an explicit lottery
over its two strict orders.

The reference matching $M$ assigns $a_v$ to $h_v$ and $b_v$ to $k_v$,
and leaves all edge applicants unmatched. It occupies every house.
For a realized profile, let $S$ contain the vertices whose variable
applicants rank $k_v$ first. The reference rank vector is
$(2n-|S|,|S|)$ when $n>0$.

\begin{lemma}\label{lem:criterion}
The reference matching is rank-maximal exactly when every edge
applicant's first choice is $h_v$ with $v\notin S$.
\end{lemma}
\begin{proof}
If an edge applicant first chooses $h_v$ with $v\in S$, remove
$(a_v,h_v)$ from $M$ and assign that previously unmatched applicant
to $h_v$. This feasible matching gains one rank-one edge and loses
one rank-two edge, so it improves the rank vector.

Conversely, suppose all edge applicants first choose inactive homes.
Every first-choice house belongs to
\[
 \{k_v:v\in V\}\ \cup\ \{h_v:v\notin S\}.
\]
Thus any matching has at most $2n-|S|$ rank-one edges, since assigned
houses are distinct. The reference attains this bound. A competitor
attaining it has at most $2n$ assigned edges in total and hence at most
$|S|$ rank-two edges. All acceptable lists have length at most two,
so no later coordinate remains to compare.
\end{proof}

\section*{One exact query counts independent sets}
By Lemma~\ref{lem:criterion}, a successful profile must have $S$ independent.
For such an $S$, let $d(S)$ count the graph edges incident to $S$.
Each of their $t$ copies has exactly one successful order; a copy of
an edge disjoint from $S$ has two. The variable orders are fixed by $S$.
The successful-profile count $N_t$ therefore satisfies
\begin{equation}\label{eq:count}
 N_t=\sum_{S\in\IS(G)}2^{t(m-d(S))},
 \qquad \eta_t=\frac{N_t}{2^{n+tm}}.
\end{equation}
Here $\IS(G)$ includes the empty independent set.

Choose $t=n+1$ and put $B=2^{n+1}$. If $c_j$ counts the independent
sets with $d(S)=j$, then
\[
 N_t=\sum_{j=0}^{m}c_jB^{m-j},\qquad 0\leq c_j\leq2^n<B.
\]
These coefficients are exactly the base-$B$ digits of $N_t$.
One exact probability query gives $N_t$ after multiplication by
$2^{n+(n+1)m}$; summing its digits gives $|\IS(G)|$.
The constructed instance has $2n+(n+1)m$ applicants and $2n$ houses.
All integers in the extraction have polynomial bit length, so construction
and digit extraction take polynomial time. Independent-set counting is
$\sharpP$-complete under polynomial-time Turing reductions
\cite[\S1, problem 2; \S2, reduction 2]{provan}.
This proves the exact lower bound for both encodings.

For the threshold lower bound, put $D=2^{n+(n+1)m}$.
Queries at $q=k/D$ determine whether $N_t\geq k$. Binary search
recovers $N_t$ in $O(\log(D+1))$ queries, after which digit extraction
recovers the independent-set count. A $\sharpP$ oracle suffices to decide
every $\PP$ language, so the threshold problem is $\PP$-hard under
the stated Turing reductions. This does not assert many-one hardness.

Isolated vertices need no modification: both variable orders contribute
to the corresponding coefficient. When $m=0$, $N_t=2^n$ and $\eta_t=1$,
and the single digit counts all subsets. The empty graph gives probability
one and one independent set.

\section*{Verification and upper bounds}
For a deterministic profile with $a$ applicants and maximum list length
$R>0$, assign an edge of rank $r$ the integer weight
$(a+1)^{R-r}$. The first differing rank coordinate dominates all later
coordinates: their total possible loss is at most
\[
 a\sum_{r=j+1}^{R}(a+1)^{R-r}=(a+1)^{R-j}-1.
\]
Maximum-weight bipartite matching, permitting unmatched vertices, therefore
verifies rank maximality in polynomial bit time. The case $R=0$ is immediate.

For rational independent lotteries, choose a common denominator $D_i$
for each applicant and represent each listed mass by its integer number
of tickets. Guess one binary ticket per applicant, reject out-of-range
tickets, and accept precisely the profiles passing the verifier. The
accepting-branch count is $N=D\eta(M)$, where $D=\prod_iD_i$ is known
and has polynomial bit length. For uniform refinements, a ticket indexes
a permutation in each tie, giving $D=\prod_T|T|!$ and the same argument.
Consequently exact evaluation belongs to $\FP^{\sharpP}$. For
$q=u/v\in(0,1]$ with $v>0$, comparison is equivalent to
$vN>uD-1$. The count $vN$ remains in $\sharpP$: add a bounded ticket
guess with $v$ choices. A $\sharpP$ count exceeding a polynomial-time
computable integer is a $\PP$ predicate
\cite[Definition 2.3]{fenner}; thresholds outside $(0,1]$ are immediate.
For explicitly listed joint support,
run the verifier on each profile and add the successful masses.
This proves all upper bounds and Theorem~\ref{thm:main}.

\paragraph{Attribution and submission scope.}
The \href{https://github.com/mehmetmars7/Erdosproblems-llm-hunter/blob/00ced55ba1223fb48603e1f15028d753997ee515/attacks/open_problems/economics/gpt6_astra_ultra/C78-464.tex}{earlier GPT 6 Astra Ultra catalogue attempt}
already gives the verifier, counting upper bounds and joint-lottery
algorithm; its two independent-model lower bounds remain unresolved.
The new contribution is the gadget, \eqref{eq:count}, and digit extraction.
This submission contains the complete hand proof. A Lean companion is
developed separately and is not part of this version. The work was developed
with Codex AI assistance, without a publication-priority or external-review claim.

\clearpage
\section*{Completion Estimate}
100\%
This is the author's claimed coverage of the complete catalogue question,
not an independent verification or community acceptance rating.

\section*{Final status}
\textbf{FULL SOLUTION CLAIMED.} Exact evaluation and rational-threshold
comparison are classified for all three uncertainty encodings in the
original catalogue question. Completeness uses polynomial-time Turing
reductions. No remaining mathematical gap is claimed in this hand proof;
independent community review is pending. No existing attempt is superseded
or edited by this contribution.

\begin{thebibliography}{9}
\bibitem{survey}
K. Cechl\'arov\'a, \'{A}. Cseh and D. F. Manlove.
\emph{Selected open problems in Matching Under Preferences}.
Bulletin of the EATCS 128 (2019), \S2.1 and \S3.1.2.
\href{https://hpi.de/friedrich/docs/publications/2019/BEATCS.pdf}{Primary manuscript}.
\bibitem{provan}
J. S. Provan and M. O. Ball.
The complexity of counting cuts and of computing the probability
that a graph is connected.
\emph{SIAM Journal on Computing} 12(4) (1983), 777--788.
\href{https://doi.org/10.1137/0212053}{doi:10.1137/0212053}.
\bibitem{fenner}
S. A. Fenner, L. J. Fortnow and S. A. Kurtz.
Gap-definable counting classes.
\emph{Journal of Computer and System Sciences} 48(1) (1994), 116--148.
\href{https://cse.sc.edu/~fenner/papers/gaps.pdf}{Author manuscript}, \S2.
\end{thebibliography}
\end{document}
